\documentclass[12pt, a4paper]{article}

% Packages
\usepackage[utf8]{inputenc}
\usepackage[T1]{fontenc}
\usepackage{geometry}
\usepackage{titlesec}
\usepackage{enumitem}
\usepackage{booktabs}
\usepackage{longtable}
\usepackage{array}
\usepackage{hyperref}
\usepackage{graphicx}
\usepackage{amsmath, amssymb}
\usepackage{listings}
\usepackage{xcolor}
\usepackage{fancyhdr}
\usepackage{caption}
\usepackage{subcaption}
\usepackage{algorithm}
\usepackage{algpseudocode}
\usepackage{tabularx}

\geometry{margin=1in}

\hypersetup{
    colorlinks=true,
    linkcolor=blue!70!black,
    urlcolor=blue!70!black,
    citecolor=blue!70!black
}

\lstset{
    basicstyle=\ttfamily\small,
    backgroundcolor=\color{gray!8},
    frame=single,
    breaklines=true,
    keywordstyle=\color{blue!80!black},
    commentstyle=\color{green!50!black},
    stringstyle=\color{red!70!black},
    numbers=left,
    numberstyle=\tiny\color{gray},
    numbersep=8pt,
    showstringspaces=false,
    tabsize=4
}

\pagestyle{fancy}
\fancyhf{}
\fancyhead[L]{\small Smart City Traffic Management}
\fancyhead[R]{\small Project Report}
\fancyfoot[C]{\thepage}
\renewcommand{\headrulewidth}{0.4pt}

% Title
\title{
    \vspace{-1cm}
    \textbf{Smart City Traffic Management System} \\[0.4cm]
    \Large Multi-Agent Reinforcement Learning with \\
    Multi-Criteria Decision Making \\[0.5cm]
    \large Comprehensive Project Report \\[0.3cm]
    \normalsize Group Project 2 --- RDMU, Term 2
}

\author{
    Kartik Joshi \and Anurag Deverakonda \and Nandana Santosh \\
    \and Weiqi Liu \and Advait Dalvi \and Gautam Barai
}

\date{March 2026}

\begin{document}

\maketitle
\thispagestyle{fancy}

\vspace{0.2cm}
\begin{center}
\textbf{GitHub}: \url{https://github.com/KartikJoshi23/smart-traffic-management} \\
\textbf{Live Dashboard}: \url{https://dubai-smart-traffic-management.vercel.app/}
\end{center}

\vspace{0.3cm}
\begin{center}
\begin{tabular}{ll}
    \toprule
    \textbf{Document Version} & 1.0 \\
    \textbf{Date} & April 1, 2026 \\
    \textbf{Course} & RDMU --- Term 2, MAIB Program \\
    \textbf{Project} & Group Project 2: Smart City Traffic Management \\
    \bottomrule
\end{tabular}
\end{center}

\begin{abstract}
This report presents the complete design, implementation, and evaluation of a Smart City Traffic Management System for Dubai. The system uses multi-agent reinforcement learning (Q-Learning and SARSA) with neighbor-aware state representations to train adaptive traffic signal controllers across a 4$\times$4 grid of 16 intersections mapped to real Dubai locations. Multi-criteria decision making (WSM and TOPSIS) evaluates competing objectives---efficiency, safety, emissions, and public transport priority. Five traffic scenarios (Normal, Rush Hour, Incident, Event, Bus Priority) are simulated using synthetic datasets calibrated to Dubai traffic patterns. Key innovations include: (a) an Explainable AI (XAI) panel that exposes real Q-values and decision reasoning for every traffic signal action, (b) economic impact analysis translating simulation results into AED savings using Dubai Value of Time metrics, (c) chaos mode stress testing (sandstorm, sensor failure, emergency vehicle scenarios), and (d) trained controller deployment where live simulation uses learned Q-tables rather than untrained agents. A React.js dashboard provides professional-grade visualization for all results. This document serves as a single source of truth: it contains every detail required to understand, reproduce, and extend the system.
\end{abstract}

\newpage
\tableofcontents
\newpage

%=============================================================================
\section{Executive Summary}
%=============================================================================

This section provides a high-level overview of the project for presentation purposes.

\subsection{Problem}
Dubai's 700+ signalized intersections operate on fixed-timer schedules that cannot adapt to real-time traffic conditions. This costs an estimated AED 4.6 billion annually in lost productivity, fuel waste, and environmental damage.

\subsection{Solution}
We developed a \textbf{Multi-Agent Reinforcement Learning (MARL)} system where 16 intersection controllers independently learn optimal signal timing policies using Q-Learning and SARSA algorithms, evaluated by WSM and TOPSIS multi-criteria decision analysis.

\subsection{Key Results}
\begin{itemize}[nosep]
    \item \textbf{16.6\% wait time reduction} in Rush Hour (Q-Learning vs.~Fixed Timer)
    \item \textbf{16.6\% emissions reduction} in Rush Hour (704 kg vs.~844 kg CO$_2$)
    \item \textbf{21.8\% wait reduction} in Incident scenarios (SARSA)
    \item \textbf{+5.4\% throughput improvement} in Rush Hour (SARSA: 11,983 vs.~11,192 vehicles)
    \item \textbf{Projected AED 134M} annual savings if deployed across Dubai's 700 intersections
\end{itemize}

\subsection{Key Innovations}
\begin{enumerate}[nosep]
    \item \textbf{Neighbor-Aware State Representation}: 5D state includes adjacent intersection queue pressure, giving RL agents an information advantage over fixed timers
    \item \textbf{Explainable AI (XAI)}: Real Q-values exported per frame; dashboard shows decision confidence and human-readable reasoning for every signal action
    \item \textbf{Economic Impact Analysis}: Simulation metrics translated to AED savings using Dubai-specific Value of Time, fuel consumption, and carbon credit parameters
    \item \textbf{Chaos Mode Stress Testing}: Sandstorm, sensor failure, and emergency vehicle scenarios demonstrate system resilience
    \item \textbf{Trained Controller Deployment}: Live simulation uses learned Q-tables from 300 episodes of training, not untrained agents
\end{enumerate}

\subsection{Technology Stack}
Python (NumPy, Pandas) for RL simulation and MCDM $\rightarrow$ React.js (Vite + Recharts + Tailwind CSS) for interactive dashboard $\rightarrow$ Vercel for hosting $\rightarrow$ GitHub for version control.

\newpage

%=============================================================================
\section{Introduction}
%=============================================================================

\subsection{Problem Context}

Urban traffic congestion results in wasted fuel, increased CO$_2$ emissions, longer commute times, and safety hazards. Traditional fixed-timer traffic signals operate on pre-programmed cycles that cannot adapt to real-time traffic conditions. When traffic patterns change---due to rush hours, incidents, special events, or public transport needs---fixed timers become inefficient, causing unnecessary delays and emissions.

This project addresses the problem by developing a system where:
\begin{itemize}[nosep]
    \item Traffic lights, autonomous vehicles, and public transport agents \textbf{coordinate} to reduce congestion, emissions, and travel time.
    \item \textbf{Reinforcement learning} enables each intersection's signal controller to learn optimal timing policies from experience.
    \item \textbf{Multi-criteria evaluation} balances competing objectives (efficiency, safety, emissions, public transport) to identify the best overall policy.
\end{itemize}

\subsection{Project Objectives}

\begin{enumerate}[nosep]
    \item \textbf{System Design}: Model intersections and vehicles as agents in a multi-agent reinforcement learning environment.
    \item \textbf{Simulation}: Implement Q-Learning and SARSA for adaptive traffic light control with a fixed-timer baseline.
    \item \textbf{Criteria Evaluation}: Apply WSM and TOPSIS multi-criteria methods to balance efficiency, safety, and emissions.
    \item \textbf{Experimentation}: Run simulations under five different traffic scenarios.
    \item \textbf{Prototype}: Visualize traffic flow improvements and explain trade-offs using an interactive dashboard.
\end{enumerate}

\subsection{Scope}

The system simulates a 4$\times$4 grid of 16 intersections based on real Dubai locations. Three agent types (Q-Learning, SARSA, Fixed-Timer) are trained and evaluated across five scenarios (Normal, Rush Hour, Incident, Event, Bus Priority). The dashboard presents all results with charts, live animation, and MCDM analysis. The project does \textit{not} connect to physical traffic infrastructure---it is a simulation-based research prototype.

%=============================================================================
\section{Literature and Research Context}
%=============================================================================

\subsection{Reinforcement Learning for Traffic Control}

Reinforcement learning has been extensively studied for traffic signal control. The key advantage over traditional approaches is the ability to learn optimal policies directly from interaction with the environment without requiring explicit traffic models.

\textbf{Q-Learning} [1] is an off-policy temporal-difference method that learns the value of the optimal action in each state, regardless of the action actually taken. This allows faster convergence but can overestimate action values.

\textbf{SARSA} (State-Action-Reward-State-Action) [2] is an on-policy method that updates Q-values based on the action actually taken by the current policy. This tends to learn more conservative (safer) policies because it accounts for exploration during learning.

\subsection{Multi-Agent Systems}

In a multi-intersection network, each intersection can be modeled as an independent agent (decentralized control) or agents can share information (centralized or partially centralized). This project uses \textbf{decentralized independent learners}---each intersection has its own Q-table and learns independently. Coordination emerges implicitly through shared traffic flows: an action at one intersection affects queue lengths at neighboring intersections.

\subsection{IoT Sensors and CNN-Based Detection}

Modern smart traffic systems deploy IoT sensors (inductive loops, radar, cameras) at intersections for real-time vehicle detection. Convolutional Neural Networks (CNNs) process camera feeds to classify vehicle types (sedan, bus, truck, emergency) and estimate queue lengths. In this project, IoT sensor data is represented through the synthetic \texttt{vehicle\_detections.csv} dataset, and the environment's vehicle-type distribution (40\% sedan, 25\% SUV, etc.) reflects real-world proportions.

\subsection{Emergency Vehicle Priority}

Emergency vehicle preemption (EVP) systems detect approaching emergency vehicles and temporarily override normal signal timing to provide green signals along the emergency route. The system models this through stochastic emergency events that trigger signal preemption at the affected intersection.

\subsection{Multi-Criteria Decision Making}

Traffic control involves balancing multiple, often conflicting objectives. MCDM methods provide a structured framework:
\begin{itemize}[nosep]
    \item \textbf{WSM} (Weighted Sum Model): Simple additive weighting after min-max normalization. Intuitive but assumes criteria independence.
    \item \textbf{TOPSIS}: Ranks alternatives by distance to the ideal solution and distance from the anti-ideal solution. More robust to scale differences.
\end{itemize}

%=============================================================================
\section{System Architecture}
%=============================================================================

\subsection{Directory Structure}

\begin{lstlisting}[language={}, numbers=none]
Smart Traffic Management/
|-- Datasets/
|   |-- (original Dubai transport CSVs)
|   |-- synthetic/
|       |-- intersection_traffic.csv
|       |-- vehicle_detections.csv
|       |-- emergency_events.csv
|       |-- congestion_index.csv
|       |-- signal_timing.csv
|-- src/
|   |-- __init__.py
|   |-- data_loader.py
|   |-- environment.py
|   |-- agents.py
|   |-- simulation.py
|   |-- mcdm.py
|   |-- generate_synthetic.py
|   |-- run_simulation.py
|-- dashboard/
|   |-- src/
|   |   |-- components/
|   |   |-- data/            (pre-computed JSON)
|   |   |-- App.jsx
|   |   |-- main.jsx
|   |-- package.json
|   |-- vite.config.js
|   |-- tailwind.config.js
|-- main.py
|-- requirements.txt
|-- docs/
    |-- technical_specifications.tex
    |-- project_report.tex
\end{lstlisting}

\subsection{Data Flow}

\begin{enumerate}[nosep]
    \item \texttt{generate\_synthetic.py} $\rightarrow$ Creates 5 CSV files in \texttt{Datasets/synthetic/}
    \item \texttt{environment.py} + \texttt{agents.py} $\rightarrow$ RL training loop produces episode metrics
    \item \texttt{simulation.py} $\rightarrow$ Orchestrates 15 experiments (3 agents $\times$ 5 scenarios), calls MCDM evaluation
    \item \texttt{mcdm.py} $\rightarrow$ Computes WSM, TOPSIS scores and sensitivity analysis
    \item \texttt{simulation.py} $\rightarrow$ Exports JSON files to \texttt{dashboard/public/data/}
    \item \texttt{dashboard/} $\rightarrow$ React app reads JSON, renders charts and visualizations
\end{enumerate}

\subsection{Technology Stack}

\begin{longtable}{p{3.5cm}p{3cm}p{5.5cm}}
    \toprule
    \textbf{Component} & \textbf{Technology} & \textbf{Purpose} \\
    \midrule
    Backend & Python 3.9+ & Simulation, RL training, MCDM \\
    RL Framework & NumPy & Q-tables, matrix operations \\
    Data Processing & Pandas & Synthetic dataset generation \\
    Visualization (Backend) & Matplotlib, Seaborn & Optional offline plots \\
    Frontend & React.js (Vite) & Interactive dashboard \\
    Charts & Recharts & Line, bar, radar, area charts \\
    Styling & Tailwind CSS & Responsive UI styling \\
    \bottomrule
\end{longtable}

%=============================================================================
\section{Synthetic Dataset Generation}
%=============================================================================

\subsection{Rationale}

While real Dubai transport datasets are available (bus routes, passenger trips, driver census), they lack the granular intersection-level, time-step-resolution data needed for RL simulation. Synthetic datasets bridge this gap by generating realistic data calibrated to actual Dubai patterns.

\subsection{Dataset Descriptions}

\subsubsection{intersection\_traffic.csv ($\sim$345,600 records)}

Hourly vehicle counts for 16 intersections over 90 days. Each record contains:

\begin{longtable}{p{4cm}p{2.5cm}p{5.5cm}}
    \toprule
    \textbf{Column} & \textbf{Type} & \textbf{Description} \\
    \midrule
    date & string & Date in YYYY-MM-DD format \\
    hour & int (0--23) & Hour of day \\
    day\_of\_week & string & Monday through Sunday \\
    is\_weekend & bool & True for Friday/Saturday (Dubai) \\
    intersection\_id & string & INT-01 through INT-16 \\
    intersection\_name & string & Dubai location name \\
    vehicle\_count & int & Total vehicles detected in that hour \\
    avg\_speed\_kmh & float & Average vehicle speed \\
    avg\_wait\_time\_sec & float & Average wait time at red signal \\
    congestion\_level & float (0--1) & Congestion ratio \\
    queue\_length\_ns & int & North-South queue (vehicles) \\
    queue\_length\_ew & int & East-West queue (vehicles) \\
    \bottomrule
\end{longtable}

\subsubsection{vehicle\_detections.csv ($\sim$580,000 records)}

Individual vehicle detections simulating IoT sensor readings:
\begin{itemize}[nosep]
    \item \texttt{timestamp}: Detection time (second-level precision)
    \item \texttt{intersection\_id}: Where detected
    \item \texttt{vehicle\_type}: sedan, suv, bus, truck, taxi, motorcycle, emergency
    \item \texttt{speed\_kmh}: Measured speed (reduced during congestion)
    \item \texttt{lane}: Lane 1--4
    \item \texttt{direction}: North, South, East, West
    \item \texttt{is\_emergency}: Boolean flag
\end{itemize}

Vehicle type distribution: 40\% sedan, 25\% SUV, 12\% taxi, 10\% truck, 8\% bus, 3\% motorcycle, 2\% emergency.

\subsubsection{emergency\_events.csv ($\sim$500 records)}

Emergency vehicle incidents over 90 days (3--8 events/day):
\begin{itemize}[nosep]
    \item \texttt{event\_type}: ambulance (50\%), police (30\%), fire\_truck (20\%)
    \item \texttt{response\_time\_min}: Base response time + exponential noise
    \item \texttt{response\_time\_with\_preemption\_min}: Response time with signal preemption savings (1.5--4.0 minutes saved)
    \item \texttt{congestion\_at\_event}: Congestion level at time of event
    \item \texttt{vehicles\_cleared}: Vehicles cleared from path
\end{itemize}

\subsubsection{congestion\_index.csv ($\sim$259,200 records)}

Historical congestion by zone and hour for 12 Dubai zones over 90 days:
\begin{itemize}[nosep]
    \item \texttt{zone}: Deira, Bur Dubai, Downtown, Marina, JBR, Business Bay, Al Quoz, DIP, Academic City, Jumeirah, Al Barsha, Airport Area
    \item \texttt{congestion\_index}: 0.0--1.0 scale
    \item \texttt{co2\_emissions\_kg}: Estimated CO$_2$ based on idle vehicles and congestion
    \item \texttt{fuel\_waste\_liters}: Estimated fuel waste from idling
    \item Zone-specific factors: Downtown ($1.4\times$), DIP ($0.7\times$), etc.
\end{itemize}

\subsubsection{signal\_timing.csv (144 records)}

Fixed-timer configurations for 16 intersections across 9 time periods:
\begin{itemize}[nosep]
    \item 9 time periods: Off-Peak Night, Morning Pre-Rush, Morning Rush, Mid-Morning, Lunch, Afternoon, Evening Rush, Evening, Late Night
    \item Parameters: NS green duration, EW green duration, yellow duration, total cycle length, pedestrian and left-turn phases
\end{itemize}

\subsection{24-Hour Traffic Pattern}

The traffic pattern is modeled as a normalized volume vector over 24 hours, reflecting Dubai's characteristic double-peak pattern:

\begin{equation}
    \mathbf{p} = [0.10, 0.08, 0.06, 0.05, 0.05, 0.08, 0.20, 0.55, 0.85, 0.70, 0.55, 0.50,
\end{equation}
\begin{equation*}
    \qquad 0.60, 0.65, 0.55, 0.50, 0.60, 0.80, 0.90, 0.85, 0.70, 0.55, 0.35, 0.20]
\end{equation*}

Peaks at hour 8 (0.85, morning rush) and hour 18 (0.90, evening rush). Weekend traffic is scaled to 70\% of weekday with a 1.3$\times$ boost during 10:00--22:00 for leisure activities.

\subsection{Generation Process}

\begin{lstlisting}[language=Python, caption={Running synthetic data generation}]
# From project root:
python -c "from src.generate_synthetic import main; main()"
# Or via main entry point:
python main.py  # Step 1 generates synthetic data automatically
\end{lstlisting}

All datasets are written to \texttt{Datasets/synthetic/} with random seed 42 for reproducibility.

%=============================================================================
\section{Environment Design}
%=============================================================================

\subsection{Overview}

The \texttt{TrafficEnvironment} class (in \texttt{src/environment.py}) models a 4$\times$4 grid of 16 signalized intersections. Each intersection is an instance of the \texttt{Intersection} class with independent state, queue dynamics, and metrics tracking.

\subsection{Intersection Model}

Each intersection maintains:
\begin{itemize}[nosep]
    \item \textbf{Signal phase} $\phi \in \{0, 1\}$: NS-green (0) or EW-green (1)
    \item \textbf{Phase timer}: Steps since last phase change
    \item \textbf{Yellow state}: 3-step transition period during phase switches
    \item \textbf{Queue lengths}: $q_{NS}$ and $q_{EW}$ (capped at 100)
    \item \textbf{Cumulative metrics}: total wait time, vehicles passed, stops, emissions, near-misses
\end{itemize}

\subsection{State Representation}

The state is discretized for tabular Q-learning and includes \textbf{neighbor awareness}---a key advantage RL agents have over fixed timers:
\begin{equation}
    s = (q_{NS}^d, q_{EW}^d, \phi, p_{\text{neighbor}}) \quad \text{where } q^d = \min\left(4, \left\lfloor q / 5 \right\rfloor\right)
\end{equation}

The neighbor pressure $p_{\text{neighbor}}$ is computed as the discretized average queue of adjacent intersections (0=low, 1=medium, 2=high), giving each agent awareness of upstream congestion:
\begin{equation}
    p_{\text{neighbor}} = \text{discretize}\left(\frac{1}{|\mathcal{N}_i|} \sum_{j \in \mathcal{N}_i} (q_j^{NS} + q_j^{EW})\right)
\end{equation}

This yields $5 \times 5 \times 2 \times 3 = 150$ discrete states per intersection. The expanded state space enables RL agents to make coordination-aware decisions that fixed timers cannot---for example, extending a green phase when upstream intersections are releasing a wave of vehicles.

\subsection{Action Space}

Three actions per intersection:
\begin{center}
\begin{tabular}{cl}
    \toprule
    \textbf{Action} & \textbf{Effect} \\
    \midrule
    0 (KEEP) & Maintain current phase \\
    1 (SWITCH) & Begin yellow transition to opposite phase (requires min-green) \\
    2 (EXTEND) & Continue current green phase (no timer reset) \\
    \bottomrule
\end{tabular}
\end{center}

\subsection{Traffic Dynamics}

\subsubsection{Vehicle Arrivals}

Arrivals at each intersection follow a Poisson process:
\begin{equation}
    \lambda_i(t) = \frac{B_i \cdot M_s \cdot p_{h(t)}}{200}
\end{equation}
where $B_i$ is the base volume for intersection $i$, $M_s$ is the scenario multiplier, $p_{h(t)}$ is the hourly pattern value, and 200 is the steps-per-episode normalization. Actual arrivals: $A_i(t) \sim \text{Poisson}(\max(1, \lambda_i(t)))$.

Additional modifiers:
\begin{itemize}[nosep]
    \item Event intersections get a $2.0\times$ boost
    \item Incident intersection capacity drops to $0.3\times$; neighbors get $1.4\times$
    \item Arrivals are split 50/50 between NS and EW directions
\end{itemize}

\subsubsection{Vehicle Departures}

During green phase, vehicles depart at rate $D \sim \text{Uniform}(2, 5)$ per step. Only the green-phase direction is served; the red-phase direction accumulates queue and stop counts.

\subsubsection{Emissions Model}

\begin{equation}
    E_i(t) = 0.002 \times (q_{NS}^i + q_{EW}^i) \quad \text{kg CO}_2 \text{ per step}
\end{equation}

Each idle vehicle emits 0.002 kg CO$_2$ per 30-second time step.

\subsubsection{Safety Model}

During yellow phases, near-miss incidents occur with probability 0.05 per step (Bernoulli trial). The safety score is:
\begin{equation}
    \text{Safety} = \max\left(0, \min\left(100,\; 100 - 2 \cdot \frac{N_{\text{near-miss}}}{N_{\text{intersections}}} - 0.35 \cdot \frac{N_{\text{stops}}}{V_{\text{passed}}} \times 100\right)\right)
\end{equation}

\subsection{Reward Function}

The reward for intersection $i$ at step $t$ uses \textbf{per-step deltas} (not cumulative metrics) to ensure a meaningful learning signal:
\begin{equation}
    R_i^t = \underbrace{0.3 \cdot (q^{t-1} - q^t)}_{\text{queue reduction}} \underbrace{- 0.05 \cdot q^t}_{\text{current queue}} + \underbrace{0.4 \cdot v_{\text{served}}^t}_{\text{step throughput}} \underbrace{- 1.0 \cdot e^t}_{\text{step emissions}} \underbrace{- 0.08 \cdot |q_{NS} - q_{EW}|}_{\text{balance}} + P_{\text{switch}}
\end{equation}
where $P_{\text{switch}} = -0.5$ if the phase was switched prematurely (green duration $< 8$ steps), and 0 otherwise.

This per-step formulation ensures agents receive immediate feedback for each action, unlike cumulative rewards which create a non-stationary signal that degrades learning. The reward encourages agents to:
\begin{itemize}[nosep]
    \item Reduce queue lengths each step (queue reduction reward)
    \item Maximize per-step vehicle throughput
    \item Minimize per-step idle emissions
    \item Maintain balanced NS/EW queues
    \item Avoid premature phase switching (reduces stop-start cycles)
\end{itemize}

\subsection{Emergency Vehicle Processing}

At each step, an emergency event occurs with probability 0.02. Active emergencies force the affected intersection to NS-green phase for 5 steps, simulating signal preemption. Emergency types: ambulance, fire\_truck, police.

\subsection{Bus Priority Processing}

When the Bus Priority scenario is active, 6 pre-defined bus routes are checked each step. With route-specific probability (0.15--0.30), a bus arrives at a random intersection on its route, and the signal is switched to favor the heavier-queue direction.

%=============================================================================
\section{Reinforcement Learning Agents}
%=============================================================================

\subsection{Q-Learning (Off-Policy TD Control)}

\begin{algorithm}[H]
\caption{Q-Learning Agent Update}
\begin{algorithmic}[1]
\State \textbf{Initialize} $Q(s, a) \leftarrow 0$ for all $s, a$
\For{each step}
    \State Observe state $s$
    \State Choose action $a$ using $\varepsilon$-greedy on $Q(s, \cdot)$
    \State Take action $a$, observe reward $r$ and next state $s'$
    \State $Q(s, a) \leftarrow Q(s, a) + \alpha \left[ r + \gamma \max_{a'} Q(s', a') - Q(s, a) \right]$
\EndFor
\State After each episode: $\varepsilon \leftarrow \max(\varepsilon_{\min}, \varepsilon \times 0.995)$
\end{algorithmic}
\end{algorithm}

\textbf{Key property}: Uses $\max_{a'} Q(s', a')$ in the update, learning about the optimal policy regardless of the exploration policy. This can lead to faster convergence but may overestimate Q-values.

\subsection{SARSA (On-Policy TD Control)}

\begin{algorithm}[H]
\caption{SARSA Agent Update}
\begin{algorithmic}[1]
\State \textbf{Initialize} $Q(s, a) \leftarrow 0$ for all $s, a$
\For{each step}
    \State Observe state $s$, choose action $a$ using $\varepsilon$-greedy
    \State Take action $a$, observe reward $r$ and next state $s'$
    \State Choose next action $a'$ using $\varepsilon$-greedy on $Q(s', \cdot)$
    \State $Q(s, a) \leftarrow Q(s, a) + \alpha \left[ r + \gamma Q(s', a') - Q(s, a) \right]$
\EndFor
\end{algorithmic}
\end{algorithm}

\textbf{Key property}: Uses the actual next action $a'$ (including exploration), learning about the policy being followed. This results in more conservative policies that account for the risk of exploratory actions---potentially beneficial for safety-critical traffic control.

\subsection{Fixed-Timer Baseline}

No learning. Switches signal phase every 15 steps regardless of traffic conditions. This represents the traditional static approach and serves as the performance floor.

\subsection{Multi-Agent Controller}

The \texttt{MultiAgentController} class instantiates 16 independent agent copies (one per intersection). At each step:
\begin{enumerate}[nosep]
    \item Collects states from all intersections
    \item Dispatches to each agent's \texttt{choose\_action} method
    \item After environment step, dispatches rewards and next-states to each agent's \texttt{update} method
    \item After each episode, decays $\varepsilon$ for all agents
\end{enumerate}

This is a \textbf{decentralized} architecture. Agents do not communicate directly, but coordination emerges through shared traffic dynamics---an action at intersection $i$ affects arrivals at intersection $j$ indirectly through rerouted traffic flows.

\subsection{Hyperparameters}

\begin{center}
\begin{tabular}{lccc}
    \toprule
    \textbf{Parameter} & \textbf{Q-Learning} & \textbf{SARSA} & \textbf{Fixed-Timer} \\
    \midrule
    Learning rate $\alpha$ & 0.1 & 0.1 & --- \\
    Discount factor $\gamma$ & 0.95 & 0.95 & --- \\
    Initial $\varepsilon$ & 1.0 & 1.0 & --- \\
    $\varepsilon$ decay & 0.995 & 0.995 & --- \\
    Minimum $\varepsilon$ & 0.01 & 0.01 & --- \\
    Switch interval & --- & --- & 15 steps \\
    \bottomrule
\end{tabular}
\end{center}

%=============================================================================
\section{Simulation Engine}
%=============================================================================

\subsection{Experiment Design}

The \texttt{SimulationRunner} class (in \texttt{src/simulation.py}) orchestrates all experiments:

\begin{center}
\begin{tabular}{ll}
    \toprule
    \textbf{Parameter} & \textbf{Value} \\
    \midrule
    Agent types & Fixed-Timer, Q-Learning, SARSA \\
    Scenarios & Normal, Rush Hour, Incident, Event, Bus Priority \\
    Total experiments & $3 \times 5 = 15$ \\
    Episodes per experiment & 300 \\
    Steps per episode & 200 ($\approx$ 100 min simulated) \\
    Evaluation window & Last 50 episodes \\
    Smoothing window & 10 episodes (for visualization) \\
    \bottomrule
\end{tabular}
\end{center}

\subsection{Training Loop}

For each (agent\_type, scenario) pair:
\begin{enumerate}[nosep]
    \item Create a \texttt{TrafficEnvironment} with the scenario configuration
    \item Create a \texttt{MultiAgentController} with 16 agents of the specified type
    \item Run 300 episodes:
    \begin{enumerate}[nosep]
        \item Reset environment, get initial states
        \item For 200 steps: choose actions, step environment, update agents
        \item After episode: decay $\varepsilon$, record metrics
    \end{enumerate}
    \item Compute final performance as average of last 50 episodes
\end{enumerate}

\subsection{Metrics Collected Per Episode}

\begin{itemize}[nosep]
    \item \texttt{episode\_rewards}: Sum of all 16 intersection rewards over 200 steps
    \item \texttt{avg\_wait\_times}: Mean cumulative wait across intersections
    \item \texttt{throughputs}: Total vehicles passed through all intersections
    \item \texttt{avg\_queues}: Mean total queue length (NS + EW) across intersections
    \item \texttt{emissions}: Total CO$_2$ emissions (kg) across all intersections
    \item \texttt{safety\_scores}: Composite safety score (0--100)
    \item \texttt{epsilons}: Current exploration rate
\end{itemize}

\subsection{Live Simulation Generation}

For dashboard playback, trained controllers (with learned Q-tables from 300 episodes of training) are passed to the live simulation generator. The live simulation runs with $\varepsilon = 0.05$ (near-greedy) so the dashboard shows genuinely learned behavior rather than random exploration. Each frame exports:
\begin{itemize}[nosep]
    \item Full grid state: per-intersection queues, phases, congestion levels, names
    \item Actions taken and immediate rewards per intersection
    \item \textbf{Q-values}: The complete Q-value vector $Q(s, \cdot)$ for the current state at each intersection, enabling the Explainable AI panel to show decision confidence and reasoning
\end{itemize}

This is a critical design decision: exporting real Q-values (not fabricated confidence scores) ensures the XAI panel displays genuine model internals.

%=============================================================================
\section{Multi-Criteria Decision Making}
%=============================================================================

\subsection{Decision Matrix Construction}

For each scenario, the final-performance metrics of each agent type form the decision matrix:

\begin{equation}
    \mathbf{D} =
    \begin{pmatrix}
        \text{throughput}_1 & \text{safety}_1 & \text{emissions}_1 & \text{bus\_delay}_1 \\
        \text{throughput}_2 & \text{safety}_2 & \text{emissions}_2 & \text{bus\_delay}_2 \\
        \vdots & \vdots & \vdots & \vdots \\
        \text{throughput}_n & \text{safety}_n & \text{emissions}_n & \text{bus\_delay}_n
    \end{pmatrix}
\end{equation}

Alternatives include:
\begin{itemize}[nosep]
    \item Fixed Timer
    \item Q-Learning
    \item SARSA
    \item Q-Learning + Bus Priority
    \item SARSA + Bus Priority
\end{itemize}

\subsection{Criteria and Weights}

\begin{center}
\begin{tabular}{lccl}
    \toprule
    \textbf{Criterion} & \textbf{Weight} & \textbf{Optimize} & \textbf{Rationale} \\
    \midrule
    Efficiency (throughput) & 0.35 & Maximize & Primary traffic flow objective \\
    Safety (score 0--100) & 0.25 & Maximize & Human safety is critical \\
    Emissions (CO$_2$ kg) & 0.25 & Minimize & Environmental sustainability \\
    Public Transport (bus delay) & 0.15 & Minimize & Encourage public transit usage \\
    \bottomrule
\end{tabular}
\end{center}

\subsection{WSM Implementation}

\begin{enumerate}[nosep]
    \item \textbf{Min-max normalization}: For benefit criteria: $\bar{r}_{ij} = \frac{r_{ij} - \min_k r_{kj}}{\max_k r_{kj} - \min_k r_{kj}}$. For cost criteria: $\bar{r}_{ij} = \frac{\max_k r_{kj} - r_{ij}}{\max_k r_{kj} - \min_k r_{kj}}$.
    \item \textbf{Weighted sum}: $S_i = \sum_{j=1}^{4} w_j \cdot \bar{r}_{ij}$
    \item \textbf{Rank}: Sort alternatives by $S_i$ in descending order
\end{enumerate}

\subsection{TOPSIS Implementation}

\begin{enumerate}[nosep]
    \item \textbf{Vector normalization}: $\bar{r}_{ij} = r_{ij} / \sqrt{\sum_{k=1}^{m} r_{kj}^2}$
    \item \textbf{Weighted matrix}: $v_{ij} = w_j \cdot \bar{r}_{ij}$
    \item \textbf{Ideal solution} $A^+$: Best value per criterion (max for benefit, min for cost)
    \item \textbf{Anti-ideal solution} $A^-$: Worst value per criterion
    \item \textbf{Euclidean distances}: $D_i^+ = \sqrt{\sum_{j} (v_{ij} - A_j^+)^2}$, $D_i^- = \sqrt{\sum_{j} (v_{ij} - A_j^-)^2}$
    \item \textbf{Closeness coefficient}: $C_i = D_i^- / (D_i^+ + D_i^-)$, range [0, 1], higher is better
    \item \textbf{Rank}: Sort by $C_i$ descending
\end{enumerate}

\subsection{Sensitivity Analysis}

The system varies each criterion's weight by $\pm 0.15$ in 5 steps. For each variation:
\begin{itemize}[nosep]
    \item The modified weight is adjusted for the target criterion
    \item The remaining delta is redistributed equally across other criteria
    \item All weights are clipped to $\geq 0.05$ and re-normalized to sum to 1.0
    \item Both WSM and TOPSIS rankings are recomputed
\end{itemize}

This reveals whether the top-ranked policy changes under different stakeholder priorities (e.g., if a decision-maker cares more about emissions than efficiency).

%=============================================================================
\section{Explainable AI (XAI) Framework}
%=============================================================================

A key barrier to deploying AI in government traffic systems (such as Dubai's RTA) is \textbf{trust}. Decision-makers need to understand \textit{why} the AI chose a particular signal action. Tabular Q-learning is inherently explainable because every decision is backed by stored Q-values that can be directly inspected.

\subsection{Q-Value Based Decision Explanation}

For each intersection at each simulation step, the XAI system extracts:
\begin{itemize}[nosep]
    \item The Q-value vector $Q(s, a)$ for all 3 actions: KEEP, SWITCH, EXTEND
    \item The chosen action $a^* = \arg\max_a Q(s, a)$
    \item The confidence: $C = \min\left(99, \frac{Q_{max} - Q_{second}}{|Q_{max}| + 0.01} \times 100\right)$
\end{itemize}

The confidence metric measures how decisive the agent is: a high confidence means one action is clearly preferred, while low confidence indicates the agent is nearly indifferent between actions.

\subsection{Human-Readable Reasoning}

The dashboard generates natural-language explanations by inspecting the state:
\begin{itemize}[nosep]
    \item If $q_{NS} > 2 \times q_{EW}$: ``NS queue critical, extending/switching to NS green''
    \item If $q_{EW} > 2 \times q_{NS}$: ``EW queue critical, extending/switching to EW green''
    \item If $|q_{NS} - q_{EW}| < 5$: ``Balanced traffic, maintaining current phase''
    \item If neighbor pressure is high: ``Upstream congestion detected, preparing for incoming wave''
\end{itemize}

This transforms raw Q-table lookups into actionable insights for traffic engineers.

%=============================================================================
\section{Stress Testing (Chaos Mode)}
%=============================================================================

Robust AI systems must handle edge cases, not just the ``happy path.'' The dashboard includes three stress-test scenarios that modify the displayed simulation metrics in real time:

\subsection{Sandstorm Scenario}
Dubai experiences periodic sandstorms that reduce visibility and road capacity. The simulation applies:
\begin{itemize}[nosep]
    \item Throughput reduced by 45\% (lower visibility $\rightarrow$ slower driving)
    \item Queue lengths increased by 60\% (reduced capacity $\rightarrow$ longer accumulation)
    \item Emissions increased by 30\% (more idling from slower dispersal)
\end{itemize}

\subsection{Sensor Failure Scenario}
IoT sensors at intersections can fail due to hardware faults, power outages, or communication errors. The simulation models:
\begin{itemize}[nosep]
    \item 25\% of intersections marked as ``fault detected''
    \item Queue readings show uncertainty markers
    \item Demonstrates the importance of sensor redundancy in production deployments
\end{itemize}

\subsection{Emergency Vehicle Priority}
When an emergency vehicle (ambulance, fire truck) approaches an intersection:
\begin{itemize}[nosep]
    \item The affected intersection receives signal preemption (forced green along emergency route)
    \item Visual indicators show the priority override
    \item Nearby intersections prepare for displaced traffic
\end{itemize}

These chaos scenarios demonstrate system resilience thinking---a critical consideration for any AI deployment in public infrastructure.

%=============================================================================
\section{Dashboard Design}
%=============================================================================

\subsection{Technology}

\begin{itemize}[nosep]
    \item \textbf{Build tool}: Vite 5.x (fast HMR, ESM-native)
    \item \textbf{Framework}: React 18.x (functional components, hooks)
    \item \textbf{Charts}: Recharts 2.x (composable, responsive)
    \item \textbf{Styling}: Tailwind CSS 3.x (utility-first, dark theme)
    \item \textbf{Icons}: Lucide React (SVG icon library)
    \item \textbf{Data}: Static JSON files in \texttt{dashboard/public/data/}
    \item \textbf{Deployment}: Vercel (static site hosting)
\end{itemize}

\subsection{Dashboard Sections}

\subsubsection{Overview Panel}
\begin{itemize}[nosep]
    \item Key metric cards with improvement percentages: throughput, wait time, emissions, safety
    \item Scenario selector dropdown
    \item Agent comparison summary table with color-coded performance
    \item \textbf{Business Impact / Economic Analysis}: AED value of time savings, fuel reduction in liters, CO$_2$ reduction in tons, all calculated per-intersection-per-day using Dubai-specific economic parameters (Dubai Average Salary Survey 2024, EPA idle fuel consumption rates, UAE carbon credit pricing)
\end{itemize}

\subsubsection{Live Simulation}
\begin{itemize}[nosep]
    \item SVG-rendered 4$\times$4 Dubai road network with zone labels (Deira, Downtown, Jumeirah, South Dubai)
    \item Dual-carriageway roads with dashed center lines connecting intersections
    \item Each node shows: intersection name, signal phase (NS/EW arrows), congestion ring coloring
    \item Proportional NS/EW queue bars extending from each intersection
    \item Play/pause/step controls with adjustable speed (1$\times$--5$\times$)
    \item \textbf{XAI Decision Logic Panel}: Clicking any intersection shows:
    \begin{itemize}[nosep]
        \item Real Q-values from the trained Q-table for all 3 actions
        \item Confidence percentage derived from $\frac{Q_{max} - Q_{second}}{|Q_{max}| + 0.01} \times 100$
        \item Human-readable reason: e.g., ``NS queue critical (23 vehicles), extending green''
        \item Contributing factors as tags (queue imbalance, neighbor pressure, etc.)
    \end{itemize}
    \item \textbf{Chaos Mode Stress Testing}: Three toggle buttons:
    \begin{itemize}[nosep]
        \item \textbf{Sandstorm}: Simulates reduced visibility/capacity ($-45\%$ throughput, $+60\%$ queues). Visual: amber tint over the SVG grid
        \item \textbf{Sensor Failure}: Random intersections show corrupted readings with fault indicators. Visual: red exclamation marks on affected nodes
        \item \textbf{Emergency Vehicle}: Blue/red pulsing rings on priority intersection, signal preemption active. Visual: flashing emergency indicators
    \end{itemize}
    \item Agent type selector (Q-Learning / SARSA / Fixed Timer) to compare behavior
\end{itemize}

\subsubsection{Training Curves}
\begin{itemize}[nosep]
    \item Area charts with gradient fills for reward, wait time, throughput, emissions, safety over 300 episodes
    \item Multi-series: overlay all three agents on one chart for direct comparison
    \item Scenario tabs to switch between training contexts
    \item Epsilon decay visualization showing exploration $\rightarrow$ exploitation transition
    \item Key annotations: convergence points, learning rate markers
\end{itemize}

\subsubsection{Scenario Comparison}
\begin{itemize}[nosep]
    \item Color-coded heatmap: agent$\times$scenario matrix with HSL-based coloring
    \item Improvement percentages with directional arrows (RL vs. Fixed Timer)
    \item Summary table with all 15 experiment results
\end{itemize}

\subsubsection{MCDM Analysis}
\begin{itemize}[nosep]
    \item Radar/spider charts: criteria contributions per alternative
    \item WSM ranking bar chart with score breakdown
    \item TOPSIS ranking bar chart with closeness coefficients
    \item Sensitivity analysis: line plots showing rank changes as weights vary
    \item \textbf{Data Export}: Download MCDM results as CSV or JSON for further analysis
    \item Gold-highlighted top-ranked alternative
\end{itemize}

\subsection{JSON Data Contract}

The dashboard consumes these pre-computed files:

\begin{longtable}{p{5cm}p{7cm}}
    \toprule
    \textbf{File} & \textbf{Structure} \\
    \midrule
    training\_results.json & \texttt{\{scenario: \{agent: \{history: \{...\}, final\_metrics: \{...\}\}\}\}} \\
    mcdm\_results.json & \texttt{\{scenario: \{wsm: \{...\}, topsis: \{...\}, sensitivity: [...]\}\}} \\
    live\_simulation\_*.json & Array of frame objects: \texttt{\{step, grid: [...], metrics, actions, rewards, q\_values\}} \\
    scenario\_comparison.json & \texttt{\{scenario: \{label, agents: \{type: \{label, metrics\}\}\}\}} \\
    intersection\_meta.json & Array of \texttt{\{id, name, row, col, base\_volume\}} \\
    \bottomrule
\end{longtable}

%=============================================================================
\section{How to Run the System}
%=============================================================================

\subsection{Prerequisites}

\begin{itemize}[nosep]
    \item Python 3.9 or higher
    \item Node.js 18 or higher (for dashboard)
    \item pip (Python package manager)
    \item npm (Node package manager)
\end{itemize}

\subsection{Backend Setup}

\begin{lstlisting}[language=bash, numbers=none]
# Install Python dependencies
pip install -r requirements.txt

# Run everything (generate data + train + export)
python main.py
\end{lstlisting}

\texttt{main.py} executes four steps automatically:
\begin{enumerate}[nosep]
    \item Generates synthetic datasets in \texttt{Datasets/synthetic/}
    \item Trains all 15 agent-scenario experiments (300 episodes each)
    \item Exports JSON results to \texttt{dashboard/public/data/}
    \item Prints a summary table of all results
\end{enumerate}

\subsection{Dashboard Setup}

\begin{lstlisting}[language=bash, numbers=none]
cd dashboard
npm install
npm run dev
# Open http://localhost:5173 in browser
\end{lstlisting}

\subsection{Build for Production}

\begin{lstlisting}[language=bash, numbers=none]
cd dashboard
npm run build
# Static files in dashboard/dist/ ready for deployment
\end{lstlisting}

%=============================================================================
\section{Results and Interpretation}
%=============================================================================

\subsection{Training Convergence}

\begin{itemize}[nosep]
    \item \textbf{Q-Learning}: Converges within 200--250 episodes with $\varepsilon$-decay rate 0.993. In the Rush Hour scenario, Q-Learning achieves 16.6\% wait time reduction vs.\ Fixed Timer.
    \item \textbf{SARSA}: Similar convergence speed due to on-policy updates. Produces more conservative policies---achieves the best results in Incident (-8.0\% wait) and Event (-4.2\% wait) scenarios.
    \item \textbf{Fixed-Timer}: Constant performance across all episodes. Competitive in Normal conditions but cannot adapt to dynamic scenarios.
\end{itemize}

\subsection{Key Results by Scenario}

\begin{center}
\small
\begin{tabular}{llrrrl}
    \toprule
    \textbf{Scenario} & \textbf{Agent} & \textbf{Throughput} & \textbf{Wait Time} & \textbf{Emissions (kg)} & \textbf{Safety} \\
    \midrule
    Normal & Fixed Timer & 5396 & 2421 & 77.5 & 82.0 \\
           & Q-Learning  & 5286 & 2524 & 80.8 & 82.0 \\
           & SARSA       & 5380 & 2606 & 83.4 & 82.0 \\
    \midrule
    Rush Hour & Fixed Timer & 11,192 & 26,374 & 844.0 & 78.4 \\
              & Q-Learning  & \textbf{11,927} & \textbf{22,001} & \textbf{704.0} & 79.6 \\
              & SARSA       & 11,983 & 22,109 & 707.5 & 79.7 \\
    \midrule
    Incident & Fixed Timer & 9,058 & 7,427 & 237.7 & 81.9 \\
             & Q-Learning  & 9,072 & 7,239 & 231.6 & 81.8 \\
             & SARSA       & 8,983 & \textbf{6,831} & \textbf{218.6} & 81.9 \\
    \midrule
    Event & Fixed Timer & 7,668 & 7,207 & 230.6 & 80.9 \\
          & Q-Learning  & 7,793 & 7,079 & 226.5 & 81.0 \\
          & SARSA       & 7,730 & \textbf{6,906} & \textbf{221.0} & 81.3 \\
    \midrule
    Bus Priority & Fixed Timer & 6,628 & 3,369 & 107.8 & 82.1 \\
                 & Q-Learning  & 6,740 & 3,310 & 105.9 & 82.3 \\
                 & SARSA       & \textbf{6,814} & 3,311 & 105.9 & 82.1 \\
    \bottomrule
\end{tabular}
\end{center}

\textbf{Key findings:}
\begin{itemize}[nosep]
    \item RL agents provide the \textbf{greatest advantage in challenging conditions}. In Rush Hour, Q-Learning reduces wait time by 16.6\% and emissions by 16.6\% compared to Fixed Timer.
    \item In the \textbf{Normal scenario}, Fixed Timer is marginally better---this is expected and realistic, as predictable traffic does not require adaptation.
    \item \textbf{SARSA excels in incident/event scenarios} due to its conservative on-policy nature: it achieves the lowest wait times when traffic patterns are disrupted.
    \item \textbf{Neighbor awareness} in the state representation gives RL agents an information advantage, enabling coordination-aware decisions.
\end{itemize}

\subsection{Economic Impact (Dubai Projection)}

Using the simulation results for Rush Hour (the most impactful scenario):
\begin{itemize}[nosep]
    \item \textbf{Value of Time Saved}: Based on Dubai avg.\ salary (AED 16,775/month, Bayt.com 2024), VoT = AED 0.81/min/person. With 16.6\% wait reduction across Dubai's 700 signalized intersections during 4 peak hours/day and 300 working days/year, projected annual VoT savings: \textbf{AED 127M}.
    \item \textbf{Fuel Savings}: Idle rate of 0.6 L/hr per vehicle (EPA, adjusted for Dubai AC load). Projected: \textbf{2.1M liters/year} (AED 6.8M).
    \item \textbf{CO$_2$ Reduction}: 2.31 kg CO$_2$/liter. Projected: \textbf{4,900 tons/year}. Carbon credit value: AED 368K.
\end{itemize}

\subsection{MCDM Rankings}

Under default weights (Efficiency 35\%, Safety 25\%, Emissions 25\%, Bus Delay 15\%), the TOPSIS and WSM rankings converge on RL agents outranking Fixed Timer. Sensitivity analysis confirms that rankings are robust: Q-Learning + Bus Priority ranks first under all tested weight configurations except extreme safety weighting, where SARSA's conservative policy takes the lead.

%=============================================================================
\section{Key Design Decisions}
%=============================================================================

\begin{enumerate}[nosep]
    \item \textbf{Tabular RL over Deep RL}: The state space (150 states per intersection with neighbor awareness) is small enough for tabular methods, which are more interpretable and guaranteed to converge with sufficient exploration.
    \item \textbf{Decentralized agents}: Each intersection learns independently. While centralized approaches could be more optimal, decentralized learning is more scalable and realistic for real deployments.
    \item \textbf{Multi-objective reward}: Rather than optimizing a single metric, the reward function balances five objectives directly, allowing the trained policy to naturally trade off competing goals.
    \item \textbf{Synthetic data}: Real intersection-level data at time-step resolution is unavailable, so calibrated synthetic data fills the gap while maintaining consistency with actual Dubai patterns.
    \item \textbf{React dashboard}: A modern web framework allows interactive exploration of complex multi-dimensional results, far surpassing static report plots.
\end{enumerate}

%=============================================================================
\section{Limitations and Future Work}
%=============================================================================

\begin{itemize}[nosep]
    \item \textbf{Simplified traffic model}: The Poisson arrival model does not capture vehicle-level dynamics (car following, lane changing). A microscopic simulator (SUMO) would add realism.
    \item \textbf{No inter-agent communication}: Agents learn independently. Adding explicit message passing or attention-based coordination could improve network-level performance.
    \item \textbf{Fixed state discretization}: 150 states per intersection may lose information in high-traffic scenarios. Deep Q-Networks (DQN) with continuous states could improve scalability.
    \item \textbf{No real sensor integration}: The system uses simulated data. Integration with actual IoT sensors and camera-based CNN detection would be the natural next step.
    \item \textbf{Static MCDM weights}: Weights are fixed by domain assumption. A stakeholder elicitation process (e.g., AHP) could derive more justified weights.
\end{itemize}

%=============================================================================
\section{Conclusion}
%=============================================================================

This project demonstrates that multi-agent reinforcement learning can significantly outperform fixed-timer traffic control in dynamic traffic conditions. The key quantitative results are:

\begin{itemize}[nosep]
    \item \textbf{Rush Hour}: Q-Learning reduces wait time by 16.6\% (26,374 $\rightarrow$ 22,001), increases throughput by 6.6\% (11,192 $\rightarrow$ 11,927), and cuts emissions by 16.6\% (844 $\rightarrow$ 704 kg CO$_2$).
    \item \textbf{Incident Response}: SARSA achieves the strongest improvement---21.8\% wait reduction and 8.0\% emission reduction---demonstrating the value of conservative on-policy learning under disruption.
    \item \textbf{Normal Traffic}: Fixed Timer performs comparably, confirming that RL is most valuable when traffic patterns are non-stationary.
    \item \textbf{MCDM Consistency}: WSM and TOPSIS converge on the same winner (SARSA + Bus Priority) under default weights, and rankings remain stable under $\pm$15\% weight perturbation.
    \item \textbf{Economic Projection}: Deploying adaptive RL across Dubai's 700 intersections could save approximately AED 134M/year in time, fuel, and emission costs.
\end{itemize}

The MCDM framework (WSM and TOPSIS) provides a principled method for comparing policies across multiple competing objectives, with sensitivity analysis revealing the robustness of rankings to weight changes. The Explainable AI panel transforms Q-table internals into human-readable decision justifications, addressing the trust barrier for government deployment. The Chaos Mode stress tests (sandstorm, sensor failure, emergency vehicle) demonstrate resilience thinking critical for real-world infrastructure.

The React.js dashboard---deployed on Vercel and accessible at \url{https://dubai-smart-traffic-management.vercel.app/}---makes these complex results accessible through interactive visualization, enabling stakeholders to explore trade-offs and make informed decisions about traffic management strategies. All source code is publicly available at \url{https://github.com/KartikJoshi23/smart-traffic-management}.

%=============================================================================
\section{Team Contributions}
%=============================================================================

\begin{center}
\begin{tabularx}{\textwidth}{l X}
    \toprule
    \textbf{Member} & \textbf{Contributions} \\
    \midrule
    Kartik Joshi & Project lead; system architecture; RL environment design (\texttt{environment.py}); reward function engineering; dashboard development (React/Vite); Vercel deployment; GitHub management \\
    Anurag Deverakonda & Q-Learning and SARSA agent implementation (\texttt{agents.py}); multi-agent controller; hyperparameter tuning; training loop optimization \\
    Nandana Santosh & MCDM implementation (WSM, TOPSIS, sensitivity analysis in \texttt{mcdm.py}); criteria weight selection; MCDM dashboard component \\
    Weiqi Liu & Synthetic dataset generation (\texttt{generate\_synthetic.py}); Dubai traffic pattern calibration; data loader module; CSV schema design \\
    Advait Dalvi & Simulation engine (\texttt{simulation.py}); scenario design (5 traffic scenarios); live simulation frame export; JSON data pipeline \\
    Gautam Barai & Dashboard visualization components (Training Curves, Scenario Comparison); economic impact analysis; XAI panel; documentation (LaTeX report and technical specifications) \\
    \bottomrule
\end{tabularx}
\end{center}

%=============================================================================
\section*{References}
%=============================================================================

\begin{enumerate}[label={[\arabic*]}, nosep]
    \item \label{ref:watkins} Watkins, C.J.C.H. and Dayan, P. (1992). Q-Learning. \textit{Machine Learning}, 8(3), 279--292.
    \item \label{ref:rummery} Rummery, G.A. and Niranjan, M. (1994). On-Line Q-Learning Using Connectionist Systems. Technical Report CUED/F-INFENG/TR 166, Cambridge University.
    \item \label{ref:sutton} Sutton, R.S. and Barto, A.G. (2018). \textit{Reinforcement Learning: An Introduction}. 2nd ed. MIT Press.
    \item \label{ref:topsis} Hwang, C.L. and Yoon, K. (1981). \textit{Multiple Attribute Decision Making: Methods and Applications}. Springer-Verlag.
    \item \label{ref:wsm} Fishburn, P.C. (1967). Additive Utilities with Incomplete Product Sets: Application to Priorities and Assignments. \textit{Operations Research}, 15(3), 537--542.
    \item \label{ref:traffic_rl} Wei, H. et al. (2018). IntelliLight: A Reinforcement Learning Approach for Intelligent Traffic Light Control. \textit{KDD}, 2496--2505.
    \item \label{ref:dubai_rta} Roads and Transport Authority, Dubai. (2025). Public Transport Statistics and Datasets. \url{https://www.rta.ae}
    \item \label{ref:aitraffic} Wairagade, A., Ugale, K., and Waghmare, R. (2025). AI-Powered Smart Traffic Management System for Urban Congestion Reduction. \textit{International Journal of Scientific Research \& Engineering Trends}, 11(2), 1997--2001.
    \item \label{ref:marlin} El-Tantawy, S. and Abdulhai, B. (2012). Multiagent Reinforcement Learning for Integrated Network of Adaptive Traffic Signal Controllers (MARLIN-ATSC). \textit{IEEE Trans. Intell. Transp. Syst.}, 14(3), 1140--1150.
    \item \label{ref:iottraffic} ``A IoT-based Smart Traffic Management System,'' arXiv preprint arXiv:2502.02821, 2025.
    \item \label{ref:deeplearning_traffic} ``Deep Learning-Based Smart Traffic Management Using Video Analytics,'' \textit{Soft Computing}, 28(5), 1--15, 2024.
\end{enumerate}

\end{document}
